% New standalone audit/partial increment. No prior source is modified.
\documentclass[11pt]{article}
\usepackage[T1]{fontenc}\usepackage[utf8]{inputenc}\usepackage{lmodern}
\usepackage[margin=.87in]{geometry}
\usepackage{amsmath,amssymb,amsthm,mathtools,microtype,array,longtable,booktabs,xurl}
\usepackage[hidelinks]{hyperref}
\setlength{\emergencystretch}{3em}\allowdisplaybreaks[2]
\numberwithin{equation}{section}
\newtheorem{proposition}{Proposition}[section]
\newtheorem{lemma}[proposition]{Lemma}
\theoremstyle{remark}\newtheorem{remark}[proposition]{Remark}
\newcommand{\dd}{\,d}\newcommand{\E}{\mathbb E}
\newcommand{\KK}{\mathcal K}\newcommand{\BB}{\mathcal B}
\newcommand{\Gp}{G_+}\newcommand{\Clk}{\mathcal C_q}
\newcommand{\Kfun}{\mathsf K}\newcommand{\Ffun}{\mathsf F}
\newcommand{\eps}{\varepsilon}
\newcommand{\src}[2]{\textup{#1} \texttt{#2}}
\title{Selection-stage interface inside the weak passage:\\
full hypothesis audit and the unmatched weak-return condition}
\author{Standalone analytical increment for independent review}
\date{Version 1 --- 28 October 2026}
\begin{document}\maketitle
\begin{abstract}
RB's completed passage is used only through its certified S4 handoff.
We inventory TAILREP's perturbed selection hypotheses and the WEAK
results underlying them. The normalized core first moment, bounded
strong moments, leading selected-branch positivity and a quantitative
weak-coordinate size bound are available. The original weak return
condition is not supplied: BULK controls absolute Cartesian ratios,
whereas WEAK requires nonnegative weak inputs and a strict inward
boundary margin. We exhibit the exact finite-noise boundary inequality
that must be proved and the order-one uncertainty left by the residual
norm estimate. We stop at this unmatched contract. No selection theorem,
logistic phase match, or continuation beyond $t_b$ is asserted.
\end{abstract}

\section{Scope, notation and the quantities that must not be identified}
\label{an03siv1:scope}
Fix $\kappa=10$, $\rho\in[13/20,7/10]$, $\lambda=\rho^2$,
$Q=\lambda+q^2$ and $\ell=\log(1/q)$. All passage facts below
come from RB's S4 handoff, with the source statement named there
when its precise formula is needed. In particular, the original flow
is certified on $[t_0,t_b]$, $0<c\le c_b$, and
\begin{gather}
 H_0\asymp A_q^{-\lambda}q^k,\quad k=10(1-\lambda),\quad
 H(t_b)\asymp c_b,\quad 0<c_A\le A_q\le C_A,\notag\\
 \log\frac{H(t)}{H(s)}=Q\int_s^t(1-c)+\delta(s,t),\qquad
 |\delta(s,t)|\le C(q^{1/8}+q^2\ell)=o(1),
 \label{an03siv1:handoffclock}\\
 0\le\int_{t_0}^{t_b}c
 \le\frac{10\mathrm e E_g}{9Q(1-c_b)}+o(1),\quad
 E_g=2c_b,\quad t_b-t_0\le20\ell.
 \label{an03siv1:handoffc}
\end{gather}
We retain the fixed original cohorts $\KK,\Gp,\Sigma,\BB$,
with $\Clk\subset\Gp$, $\mathcal D=\varnothing$ and $H=H_{\Clk}$.
The bounds on residuals, core geometry, secondary energy and outside
mass are those already handed off by RB; they are not reconstructed
from conditional pre-RB passage theorems.

The chart convention is important. P \texttt{pa:charts} and BULK
\texttt{weakJ} use
\begin{equation}
 a_w=(\sin(qu),\cos(qu)),\qquad
 b_w=(-\sin(qv),\cos(qv)),\qquad
 u=\frac{\pi/2-\theta}{q},\quad
 v=\frac{\psi-\pi/2}{q}.
 \label{an03siv1:weakcoords}
\end{equation}
Thus the output coordinate has the opposite angular convention from
the input coordinate. In particular $u\ge0$ means $U_1\ge0$, not
merely $U_2>0$.

\begin{center}
\begin{tabular}{p{.21\textwidth}p{.34\textwidth}p{.34\textwidth}}
\toprule
Notation & WEAK/TAILREP meaning & Relation to passage quantities\\\midrule
$M$ & Total strong-family squared-radius mass & Candidate is actual
$\mu_{\KK}$ if the complement is charged separately. It is not
$M_{\rm res}$ or the fixed upper cap $M_K=2$.\\[3pt]
$M=1$ & Invariant surface of the defined leading equations &
CFA's capped mass $\widetilde M=\min(\mu_{\KK},1)$ is not identically
one. Normalizing a measure does not erase its field defect.\\[3pt]
$N$ & Weak-family mass with exact law $N'=\lambda N(1-N)$ &
Neither the signed response $c$, nor $H$, nor the original radial mass
$N_C=\int_Cm\dd\lambda_0$ has this law by definition.\\[3pt]
$Y$ & Strong donor moment $\int y\dd\mu_s$ & Not LCRC's transverse
guard variable $Y$ or its bound $K_Y$.\\[3pt]
$B_\rho$ & A bound for $|\rho K_s|$ in WEAK forcing & Not the
passage residual integral $\int r\le B$.\\[3pt]
$M_K$ in TAILREP & Actual mass of the charted core & Not the
order-one passage mass cap with the same letters.\\
\bottomrule
\end{tabular}
\end{center}

\section{Every hypothesis of the selection mechanism}
\label{an03siv1:inventory}
The statements being audited are exact statements for their stated
systems. Their hypotheses cannot be replaced by small residual norms
without a matching calculation.

\subsection*{The underlying characteristic and moment structure}
WEAK works with fixed strong and weak labels, positive radial measures,
$\mu_s(\mathbb R^2)=1$, and $\mu_w=N\eta_w$, where $\eta_w$ is a
probability with \emph{fixed normalized label weights}. Its class has
bounded support on each compact time interval (the bound may diverge
as $q\to0$), sufficient moments and differentiable characteristic
integrals. The exact leading weak equations are
\begin{align}
 2u'&=\phi(u)-\mathcal W(u)-Y\Phi(u)-\lambda(1-N)u,
 &\mathcal W(u)&=N\E_{\eta_w}\Ffun(u,\widetilde u),\notag\\
 2v'&=\rho K_s-\lambda V-\lambda(1-N)v,
 &V&=N\E_{\eta_w}v,\quad N'=\lambda N(1-N).
 \label{an03siv1:leading}
\end{align}
Here $K_s=\int\Kfun(\rho x)\dd\mu_s$ and
$K_w=N\E_{\eta_w}\Kfun(u)$. Fixed original labels are available;
constant normalized original weights and \eqref{an03siv1:leading}
are not original-flow identities. RB supplies fixed-distortion primary
transport and an $o(1)$ clock defect, not equality of radial reaction
rates. The normalized-weight and field discrepancies must be paid in
TAILREP's errors, rather than discarded.

\subsection*{WEAK forcing, concentration and population-to-orbit comparison}
On a leading interval $[t_s,t_n]$, WEAK \texttt{forcing} requires
$0<N_0<n_b\le1/2$, $N(t_n)=n_b$, \eqref{an03siv1:leading},
$Y\ge0$, and $|\rho K_s|\le B_\rho$. Its finite incoming moments are
\[
 U_0^+=\E(u(t_s))_+,\quad V_0^{\rm abs}=\E|v(t_s)|,\quad
 N_*=N_0(U_0^++3V_0^{\rm abs})^2,
\]
and its constants are exactly
\begin{equation}
 C_u=\phi(0)+\frac{\phi(0)}{\lambda(1-n_b)},\qquad
 C_w^{\rm force}=C_u+B_\rho/\lambda.
 \label{an03siv1:forcingconstants}
\end{equation}
It proves $K_w\le C_uN+U_0^+\sqrt{N_0N}$ and its integrated
forcing bounds \emph{in that leading system}.

WEAK \texttt{concentration} additionally requires
\begin{equation}
 u(t_s,b)\ge0\quad\hbox{labelwise},\qquad
 0\le Y\le Y_{\max}<2\phi(0),\qquad
 \mathfrak m:=\phi(0)-Y_{\max}/2-C_un_b-U_0^+\sqrt{N_0n_b}>0.
 \label{an03siv1:returnhyp}
\end{equation}
The strict margin and the exact leading equations give the invariant
region $u\ge0$ and the multiplier $\sqrt{N(s)/N(t)}$.
Neither positivity of $z=(U_2+W_2)/2$ nor positivity of the second
coordinates is this hypothesis.

WEAK \texttt{Wineq} requires that nonnegative-input conclusion,
the same full field and fixed weight probability, and the selected
coherent branch evaluated at the \emph{same logistic mass} $N(t)$.
The selected input must also be nonnegative so that the whole segment
between a receiver and its comparator lies in the contraction region.
AN06 \texttt{regiontheorem}, for the selected leading branch and
$0<N\le1/2$, supplies $0<U_*(N)<1$, $0<V_*(N)<11/10$.
It does not supply the sign of an original population's input.
For $1\le p\le\infty$ the needed finite quantity is
\[
 W_p=\big\||u-U_*|+|v-V_*|\big\|_{L^p(\eta_w)}.
\]
The resulting weak comparison uses
$\Delta_s=|Y-Y_*|+\rho|K_s-\Kfun(\rho X_*)|$.

\subsection*{The strong local comparison and its initial margin}
WEAK's derivation of the coupled error system also requires its local
strong neighborhood, not just a fixed passage chart. Write
$a=\rho\Kfun(\rho a)$, $\alpha=\lambda\Phi(\rho a)/2$, and
\[
 \gamma=\frac{\alpha(1/2-\alpha)}{\alpha+1/2},\qquad g=\gamma/2.
\]
Choose $\delta>0$ and a small leading mass level so that
\begin{equation}
 \delta<a/8,\quad a+3\delta<2\phi(0),\quad
 |X_*(N)-a|,|Y_*(N)-a|\le\delta,
 \label{an03siv1:localbox}
\end{equation}
and the symmetric coherent strong Jacobian on
$|x-a|,|y-a|\le2\delta$, $N\le n_b$, with the selected weak
atom held fixed, is at most $-gI$. The source proves the existence
of these choices analytically. The reference must stay in
$E_s=|(c_s,d_s)-(X_*,Y_*)|<\delta$ until its guard is improved.
The donor replacement needs a \emph{pointwise} strong first moment,
not just an integrated passage error. A fixed $C_s$ must bound the
strong donor perturbations on that neighborhood; WEAK allows
$C_s=20(R_*+1)^2$ with a fixed
$R_*>2+a+|u_\rho|+a/\lambda$ bounding the selected coordinates.
For WEAK's unperturbed selection theorem the pointwise assumption is
$D_s\le d_0$, with its unit-mass strong measure, and the constants in
the forcing/return step are specifically
$Y_{\max}=a+3\delta$ and
$B_\rho=\rho\phi(0)+\lambda(a+3\delta)$.
Its moment comparison is
$\Delta_s\le(1+\lambda)(E_s+D_s)$, so the weak additive coefficient
there is $C_w^0=(1+\lambda)/2$, distinct from
$C_w^{\rm force}$. The core-only version must instead use
\eqref{an03siv1:massdeficit} and charge the complement.

\subsection*{TAILREP's exact abstract hypotheses}
TAILREP \texttt{selection} assumes nonnegative error functions satisfying
\begin{align}
 D^+E_s&\le-gE_s+C_s(d_0+NW_p),\label{an03siv1:neededstrong}\\
 D^+W_p&\le[-\lambda(1-N)/2+N/2+\ell_B(t)]W_p
                             +C_w(E_s+d_0),
 \label{an03siv1:neededweak}
\end{align}
up to the strong first exit, with $d_0\ge0$, $\ell_B\ge0$,
fixed $g,C_s,C_w>0$, exact logistic $N$, $0<N_0<n_b<1$,
$N(t_n)=n_b$, and a common finite bound
\begin{equation}
 \int_{t_s}^{t_n}\ell_B\le L_*<\infty.
 \label{an03siv1:ellbudget}
\end{equation}
For the proposed use of WEAK and AN06 we need the smaller common
range $n_b\le1/2$. The weak return-region hypotheses are
\emph{explicitly maintained separately} in TAILREP; its conclusion
does not prove them.

The constants and small-gain condition, retaining their coefficients, are
\begin{gather}
 B_*=e^{L_*}(1-n_b)^{-1/(2\lambda)},\quad
 a_0=\lambda(1-n_b)/2,\quad m_0=\min(g/2,a_0/2),\notag\\
 L_0=1+\frac{C_s(1+B_*)}{g-m_0},\quad K_0=2L_0,\quad
 \Theta=\frac{C_sB_*C_wn_b}{(a_0-m_0)(g-m_0)}\le\frac14.
 \label{an03siv1:smallgain}
\end{gather}
One also needs
\begin{equation}
 E_0=E_s(t_s),\quad A_0=\sqrt{N_0}W_p(t_s),\qquad
 K_0(E_0+d_0+A_0\sqrt{n_b})\le\delta/2.
 \label{an03siv1:initialguard}
\end{equation}
The terminal convergence conclusion needs $N_0\to0$, $A_0\to0$,
$d_0\to0$ and uniformly admissible $E_0$ and constants, including
a common $L_*$. Only then does exact logistic integration imply
$t_n-t_s\to\infty$ and the decay in \texttt{Econclusion/Wconclusion}.
The theorem neither constructs \eqref{an03siv1:neededstrong}--
\eqref{an03siv1:neededweak} nor proves the original logistic matching.

\subsection*{The mass-deficit term is part of the contract}
For a physical charted core of mass $m_K=\mu_{\KK}$, TAILREP's
elementary donor comparison is
\begin{align}
 |Y_K-Y_*|&\le D_K+m_KE_s+|m_K-1||Y_*|,\notag\\
 |K_{s,K}-\Kfun(\rho X_*)|
 &\le\rho D_K+\rho m_KE_s+|m_K-1|\Kfun(\rho X_*).
 \label{an03siv1:massdeficit}
\end{align}
Moreover P's full leading equations contain
\begin{equation}
 2u'\supset-(1-M)v\Phi(u),\qquad
 2v'\supset-(1-M)\Kfun(u).
 \label{an03siv1:masslinear}
\end{equation}
Thus unit normalization of the core leaves both comparator errors
and errors linear in large weak coordinates. Their value or coefficient
budgets must be proved on the selected late interval. CFA's
$T=0$ convention removes a complement moment from the \emph{strong}
c-form representation; it does not set \eqref{an03siv1:masslinear}
to zero.

\section{Closed parts of the audit}
\begin{proposition}[Core moment bounds available for a candidate interface]
\label{an03siv1:coremoments}
Let $\bar\mu_K=\mu|_K/\mu_K$ be the positive unit-mass core measure.
For small $q$, the RB handoff gives
\begin{equation}
 \bar D_K:=\int_K(|x-c_s|+|y-d_s|)\dd\bar\mu_K
       =O(q^{1/20}),\qquad D_K=O(q^{1/20}).
 \label{an03siv1:Dcore}
\end{equation}
Its strong moments obey, for instance,
\begin{equation}
 0<\bar Y_K\le3/5,\qquad
 0<\rho\bar K_{s,K}\le3/5.
 \label{an03siv1:normalizedmoments}
\end{equation}
For the unnormalized core one may use $0<Y_K\le5/8$ and
$|\rho K_{s,K}|\le3/5$ for sufficiently small $q$.
These are core-moment bounds, not the original weak field representation.
\end{proposition}
\begin{proof}
RB supplies $r_K=O(q^{1/20})$, $e\le3e_*/4\le3/16$,
positive core coordinates and $\mu_K\le1+o(1)$ bounded away
from zero. The distance integral is at most $2r_K$ after
normalization and at most $2\mu_Kr_K$ before it.
AN06 \texttt{residentbounds} gives $a<\phi(0)<2/5$.
Consequently every core coordinate is at most
$a+3/16+r_K<3/5$ for small $q$; its positivity is the RB core
sign check, also explicitly required in this interface's input list.
Since $\Kfun(w)\le\phi(0)+w$ for $w\ge0$,
\[
 \rho\bar K_{s,K}\le\rho\phi(0)+\lambda\bar x_K
 <\frac7{25}+\frac{147}{500}=\frac{287}{500}<\frac35.
\]
Multiplication by $\mu_K\le1+o(1)$ preserves the stated larger
unnormalized bounds. In particular $Y_{\max}=5/8$ has a positive
geometric margin $\phi(0)-Y_{\max}/2>1/16$.
The remaining incoming-moment term in \eqref{an03siv1:returnhyp}
is a separate obligation. The narrower local box
\eqref{an03siv1:localbox} is not proved by these coarse bounds.
\end{proof}

\begin{proposition}[What BULK actually gives in weak coordinates]
\label{an03siv1:bulksize}
On RB's passage the clock cohort satisfies
\begin{equation}
 |u(t,b)|+|v(t,b)|
 \le C\left[q^{-15/8}e^{-Q(1-c_b)(t-t_0)/4}+1\right].
 \label{an03siv1:size}
\end{equation}
For any fixed $h>0$ below a fixed lower bound for $H(t_b)$,
let $t_h^H$ be its first clock hit. It lies before $t_b$, and
\begin{align}
 t_h^H-t_0&=\frac{k\ell+\lambda\log A_q+\log h+O_{
m fixed}(1)}Q,
 \label{an03siv1:clockhit}\\
 |u(t_h^H,b)|+|v(t_h^H,b)|&\le C_h(1+q^{-p_b}),\quad
 p_b=\frac{15}{8}-\frac{k(1-c_b)}4,
 \quad \frac{69}{160}\le p_b\le\frac{147}{160}.
 \label{an03siv1:endpointsize}
\end{align}
These estimates give size, not sign, not bounded weak coordinates,
and not a logistic weak-mass endpoint.
\end{proposition}
\begin{proof}
Use BULK Theorem 4.1 with $\eta=1/8$. Its hypotheses are exactly
the clock entry, residual, response-range, duration and integrated
$r$ facts in RB's handoff. In its notation
\[
 x_{\rm rat}=qj/z\le C[q^{1/8}e^{-b_*(t-t_0)/2}+q^2],
 \quad b_*=Q(1-c_b)/2,\quad |a_-|\le z/3,\quad z>0.
\]
Both second coordinates are at least $2z/3$, while
$|U_1|,|W_1|\le\sqrt2j$. With the upper-half-circle lifts,
\[
 |u|\le q^{-1}|U_1/U_2|\le\frac{3\sqrt2}{2q^2}x_{\rm rat},
 \qquad
 |v|\le\frac{3\sqrt2}{2q^2}x_{\rm rat}.
\]
Here $|\arctan w|\le|w|$ was used; no sign was introduced.
This proves \eqref{an03siv1:size}.
RB has a fixed $c_H>0$ with $H(t_b)\ge c_Hc_b$ for small $q$;
choose $h<c_Hc_b/2$. Continuity and $H_0\to0$ give the first hit.
Equations \eqref{an03siv1:handoffclock}--\eqref{an03siv1:handoffc}
give \eqref{an03siv1:clockhit}. Substitution proves the exponent.
Since $51/10\le k\le231/40$ and $0<c_b\le1/4$,
\[
 \frac{15}{8}-\frac{231}{160}=\frac{69}{160},\qquad
 \frac{15}{8}-\frac{153}{160}=\frac{147}{160}.
\]
For example this certifies an endpoint chart of size
$q^{-1+1/20}$ after absorbing fixed constants, because
$147/160<19/20$. At $c_b=0$ as an algebraic limit, the upper-box
value is $p_b=3/5$; at fixed positive $c_b$ the corresponding
value is $3/5+51c_b/40$. The $c_b$ correction cannot be dropped
from a power of $q$.
\end{proof}

\begin{remark}[Why the mass-level convention needs a separate match]
\label{an03siv1:masslevel}
The condition $n_b\le c_b/2$ alone does not imply that the first
$H=n_b$ hit precedes $t_b$: RB gives fixed-factor comparability,
not $H(t_b)\ge c_b/2$. Choosing a level below $c_Hc_b/2$ resolves
this issue for $H$. It still does not identify that time with the
first hit of a leading logistic $N=n_b$ or an original radial mass
$N_C=n_b$. No $t_{n_b}$ with all these meanings is used here.
\end{remark}

\section{Exact stopping point: weak return is not an absolute-ratio bound}
\label{an03siv1:stop}
The first attempted dynamical match is the separate weak-return
hypothesis of TAILREP \texttt{selection}. The proposed use of
BULK fails to establish \eqref{an03siv1:returnhyp}.
Positive $z$, positive $U_2,W_2$, and small $qj/z$ do not imply
$U_1\ge0$. For example the balanced state
\begin{equation}
 U=W=\sqrt m(-\sin q,\cos q)
 \quad\Longrightarrow\quad
 u=-1,\ v=1,\ a_-=0,\ z>0,\quad qj/z=q\tan q=O(q^2).
 \label{an03siv1:signexample}
\end{equation}
It satisfies the Cartesian size and signgood conditions, including
their $q^2$ floor, but is outside WEAK's input return region.
This is a test of the implication between inequalities, not an
assertion that the initialized flow visits this state.

Nor does the handed-off cross-residual estimate provide the needed
strict boundary direction. At $u=0$, so $\theta=\pi/2$, the exact
angular equation gives
\begin{equation}
 u'=\frac{R_{21}(\pi/2)\cos(qv)-R_{11}(\pi/2)\sin(qv)}q.
 \label{an03siv1:exactboundary}
\end{equation}
The sign of the second term agrees with the output convention
\eqref{an03siv1:weakcoords}. In the $M=1$ leading field,
\eqref{an03siv1:returnhyp} instead gives
\[
 2u'\big|_{u=0}=\phi(0)-\mathcal W(0)-Y/2\ge\mathfrak m>0,
 \qquad \mathcal W(0)\le K_w.
\]
A sufficient original-flow version with half this strict margin is
\begin{equation}
 \boxed{\quad
 \frac{R_{21}(\pi/2,t)\cos(qv)-R_{11}(\pi/2,t)\sin(qv)}q
       \ge\frac{\mathfrak m}{4}>0
 \quad}
 \label{an03siv1:missingreturn}
\end{equation}
This is required at every potential selected-cohort boundary hit.
It must be accompanied by $u(t_s,b)\ge0$ on the entire chosen
fixed cohort, including its endpoint labels.

What the available matrix norm estimate actually gives is only
\begin{equation}
 u'\big|_{u=0}\ge-C_\times(1+M_{\rm res})-r(t)|v|.
 \label{an03siv1:normfailure}
\end{equation}
Here one may fix the coarse bound with $C_\times\ge1$.
The factor $q$ in $R_{21}=O((1+M_{\rm res})q)$ is canceled by
the scaled angular derivative. Even at $v=0$ there is an
order-one, sign-indefinite remainder. To obtain
\eqref{an03siv1:missingreturn} by bounding the error around its
leading value, the attempted norm-charge method would require
\begin{equation}
 C_\times(1+M_{\rm res})+r(t)|v|
          \le\mathfrak m/4
 \quad\text{as a pointwise boundary-error budget},
 \label{an03siv1:failedabsorption}
\end{equation}
which is not supplied by RB and does not become true by sending
$q\to0$. For this coarse choice its left side is at least one,
whereas $\mathfrak m/4\le\phi(0)/4<1/10$; this particular
absolute-error absorption fails. More basically, an unsigned bound on $R_{21}$ cannot
itself yield a positive lower bound on it. Its allowed matrix-value
test $R_{21}=-a q$, $v=0$, with any fixed admissible $a>0$,
has $u'=-a$. This is again a norm-scope test, not a counterexample
to the initialized population or to RB.

The integral $\int r\le B$ cannot supply a pointwise boundary
margin. Taking $n_b$ smaller reduces leading weak feedback but
does not make the fixed $M_{\rm res}$ norm constant vanish.
The one-sided selected-moment cancellation which could prove
\eqref{an03siv1:missingreturn} belongs to a weak-receiver field
calculation; it is not part of RB's fixed strong-chart c-form
handoff. AN06's positive $U_*$ proves the comparator sign only.

This stops the selection application. No original return hypothesis
is inserted into a conditional selection theorem. The diagnosis
does not assert that \eqref{an03siv1:missingreturn} is false for the
actual flow: it identifies the exact estimate left unproved by the
attempted BULK/norm route. No second alternative route is attempted.

\section{Remaining inventory and accuracy audit, without applying later stages}
\label{an03siv1:remaining}
The following table completes the inventory. ``Unmatched'' means
not discharged for the original population in this increment, not
that the leading source theorem is incorrect.
\begin{longtable}{>{\raggedright\arraybackslash}p{.23\textwidth}
 >{\raggedright\arraybackslash}p{.31\textwidth}
 >{\raggedright\arraybackslash}p{.36\textwidth}}
\toprule
Hypothesis & Matched part / source & Unmatched original-flow item\\\midrule
Model, sign/range, duration & RB: ratio box, $0<c\le c_b$,
$t\le t_b$, logarithmic duration & $0<N_0<n_b$ and an exact logistic
$N$ are different requirements.\\[4pt]
Characteristic regularity, finite moments, $p$ & Fixed labels, finite
mass and compact initial circle; BULK's positive second coordinates
allow the weak-angle lift for fixed $q$. Take any fixed
$1\le p\le\infty$. & No uniform-in-$q$ incoming moment or support
claim follows solely from finite support at fixed $q$.\\[4pt]
Unit strong mass and donor moment & Core normalization and
\eqref{an03siv1:Dcore}; CFA representation is valid on its strong chart.
& Pay \eqref{an03siv1:massdeficit}--\eqref{an03siv1:masslinear}; need
late sup/integrated mass-defect budgets for the selection error system.
The raw RB mass interval is not $M\equiv1$.\\[4pt]
$Y\ge0$, bounded $\rho K_s$ & Proposition \ref{an03siv1:coremoments}
gives these for core moments. & Need the original weak receiver field
to use those moments with the other cohorts charged as errors.\\[4pt]
Weak incoming sign and strict margin & Leading source's exact margin
is \eqref{an03siv1:returnhyp}; selected branch $U_*>0$ is AN06.
& Original sign at $t_s$ and \eqref{an03siv1:missingreturn}: the
stopping point.\\[4pt]
Fixed normalized weak weights & RB/BULK give fixed-cohort amplitude
transport (clock $o(1)$ defect, primary fixed distortion).
& Original radial weights, choice of a fixed comparison probability,
and their induced value/linear defects need a bound in
\eqref{an03siv1:neededweak}.\\[4pt]
Local stable strong box & WEAK proves existence of
\eqref{an03siv1:localbox}, its negative Jacobian and donor constants.
& RB's $e\le3e_*/4$ is a different fixed allowance; it does not
establish \eqref{an03siv1:initialguard} for this $\delta,K_0$.\\[4pt]
Original-to-leading fields & RB supplies a fixed strong-chart
value/derivative remainder and paid secondary/outside budgets.
& \eqref{an03siv1:neededstrong}--\eqref{an03siv1:neededweak} on the
weak receiver chart, with $d_0=o(1)$ and
\eqref{an03siv1:ellbudget}.\\[4pt]
Exact logistic phase & RB gives the exact $H$ identity with $o(1)$
defect and fixed terminal comparability. & Match $N_C,H,c$ with
an $o(1)$ logit phase; matching through the $O(1)$ hitting-time
formula is insufficient.\\[4pt]
Small gain and initial smallness & Equations
\eqref{an03siv1:smallgain}--\eqref{an03siv1:initialguard} are the exact
source conditions. & A proved common $L_*$ and actual $N_0,A_0,d_0,E_0$
are needed before asserting them for the original solution.\\
\bottomrule
\end{longtable}

\paragraph{Mass matching: the precise static statement.}
For any original fixed cohort $C$, balance gives the exact identity
\begin{equation}
 N_C=\int_Cm\dd\lambda_0=H_C+A_C+J_C.
 \label{an03siv1:massidentity}
\end{equation}
BULK \texttt{weakJ} is a bounded-chart statement. If the support
is instead $|u|,|v|\le L(q)$, $qL(q)=o(1)$, elementary
trigonometry gives
\begin{equation}
 \frac{N_C}{H_C}=1+O(q^2L(q)^2).
 \label{an03siv1:growingmass}
\end{equation}
Indeed $z^2=m[\cos(qu)+\cos(qv)]^2/4$.
For $L=q^{-1+\eps}$ the error is $O(q^{2\eps})$, not
generally $O(q^2)$. The aligned state with $u=L,v=-L$ has
$N_C/H_C=\sec^2(qL)=1+q^{2\eps}+O(q^{4\eps})$.
This is still $o(1)$ for fixed $\eps>0$; it is not by itself
a dynamical logistic phase estimate. Without such a chart, the
handed-off signgood and energy bounds only give
$1\le N_C/H_C\le10/9+J_C/H_C$.
For example, if a comparison logistic $N$ is initialized at $H(t_s)$,
then the exact clock identity gives the prospective matching relation
\begin{equation}
 \log\frac{H(t)}{N(t)}
 =q^2\int_{t_s}^t(1-c)+\lambda\int_{t_s}^t(N-c)
                         +\delta(t_s,t).
 \label{an03siv1:phaseneeded}
\end{equation}
The first and third terms vanish on the passage. Uniform vanishing
of the middle signed integral is a remaining requirement; a
sufficient condition is $\int_{t_s}^{t_n}|N-c|=o(1)$.
RB's fixed $O(c_b)$ response integral alone is not that condition.

\paragraph{What the field budget would have to say.}
The original-to-leading reduction must separate additive value errors
from coefficients multiplying the potentially large $u,v,W_p$.
The required conclusion is exactly
\eqref{an03siv1:neededstrong}--\eqref{an03siv1:ellbudget}, with
\[
 d_0\ge C\sup_{[t_s,t_n]}
   \{D_K+|1-\mu_K|+\hbox{additive normalized field errors}\}=o(1).
\]
RB pays $D_K=O(q^{1/20})$ and the outside costs with explicit
positive powers, and its secondary estimate is available.
Its fixed strong-chart remainder, however, has a small fixed
$C_0E_g$ term, and its domain is $|\theta|,|\psi|\le L_cq$.
A weak receiver near $\theta=\pi/2$ is not in that chart.
Neither its norm nor the original cross-residual bound can simply
be renamed $d_0=o(1)$ on a growing weak chart. The proposed
finite-noise coefficient $q^2(1+u^2)$ also needs its value versus
linear-error placement proved, not just a size substituted into it.

\paragraph{Late start and smallness are not declared.}
A late first hit of an actual clock level $H=q^a$, $0<a<k$, can
be located using \eqref{an03siv1:handoffclock}; that does not verify
the original return sign, the full field representation, or the
mass-defect requirements. No $t_s$ is asserted here to satisfy all
these conditions. Likewise $J_G=O(q^2)$ alone permits only an
order-one bound on a scaled weighted transverse size
$\sqrt{J_G}/q$; it is not a proof of $A_0\to0$.
These are inventory observations, not attempts to run Steps 2--6
after the stopping point.

\section{Constant order and the S4b handoff}
\label{an03siv1:handoff}
Every WPT/RB constant remains fixed in its original order. New
leading comparison constants would then be chosen from the analytic
WEAK neighborhood: $\delta,g,C_s,C_w$, followed by a \emph{proved}
common upper bound $L_*$, and a fixed $n_b$ satisfying
\eqref{an03siv1:smallgain}, the return margin and the correctly
matched endpoint constraint. Only after those choices could one
choose $q$ sufficiently small and define a suitable late start.
This is an obligation list, not an assertion that $L_*$ or the
return margin has been obtained. No preceding guard constant is
changed to force \eqref{an03siv1:failedabsorption}.

The closed outputs for S4b are:
\begin{itemize}
 \item The complete selection/forcing/concentration hypothesis list,
 with the distinct meanings of mass, response, clock and transverse
 guard variables fixed.
 \item Core first-moment replacement $O(q^{1/20})$ and bounded
 positive core donor moments, with explicit coarse constants.
 \item The absolute weak-coordinate bound \eqref{an03siv1:size} and
 the exact exponent \eqref{an03siv1:endpointsize} at a fixed actual
 clock level. This is compatible with a growing chart; it proves no sign.
 \item The corrected static growing-chart mass comparison
 \eqref{an03siv1:growingmass}, and the distinction between a first
 $H$ hit and a logistic weak-mass hit.
\end{itemize}
The first unproved dynamical contract for S4b is the original input
return: construct a fixed-cohort late start with $u(t_s,b)\ge0$,
and prove \eqref{an03siv1:missingreturn} by an actual weak-receiver
selected-moment calculation, or prove another return argument with
the same consumer conclusion. Endpoints must be included; signgood
in the second coordinate is not a replacement. The field/error,
logistic phase and initial-smallness rows then remain to be discharged.

No retained assumption is appended to a claimed selection theorem.
No S4b argument, S5 core retention, or continuation past $t_b$ is
undertaken. The obstruction is a missing implication at the prescribed
interface, not a counterexample to initialized selection.

\begin{thebibliography}{9}\small
\bibitem{RB}\path{AN03_response_barrier_v1.tex}, prefix
\texttt{an03rbv1:}; S4 handoff and the exact formulas referenced there.
\bibitem{TAILREP}\path{AN03_mass_charged_tail_and_selection_defects_v1.tex},
prefix \texttt{inc:AN03:tailrepair:v1:};
\texttt{strongerror, weakerror, selection, Econclusion, Wconclusion}.
\bibitem{WEAK}\path{AN03_weak_transit_forcing_and_orbit_selection_v1.tex},
prefix \texttt{inc:AN03:weak:v1:};
\texttt{setup, forcing, concentration, barrier, Wineq, strongineq,
smallgain, selection, guard}.
\bibitem{BULK}\path{AN03_bulk_clock_and_secondary_Q2_energy_v1.tex},
prefix \texttt{inc:AN03:bulkclock:v1:}; Theorem 4.1,
\texttt{bulkratioconclusion, weakJ}.
\bibitem{AN06}\path{AN06_learned_orbit_witnesses_v2.tex},
prefix \texttt{inc:AN06:v2:}; \texttt{residentbounds, regiontheorem}.
\bibitem{P}\path{THEOREM_A_ANALYTICAL_PROGRESS.tex};
\texttt{pa:charts, pa:leading, pa:masslaws}. Used for definitions of
the leading system, not as a finite-noise reduction theorem.
\end{thebibliography}
\end{document}
